\chapter{Acquis, débats et contribution du volume}\label{ch:contribution}

\section{Ce qui est acquis}

Sur les points suivants, le volume s’appuie sur des résultats établis par la recherche antérieure et ne prétend pas les renouveler :
\begin{itemize}
\item le cadre institutionnel de l’Église romaine au \textsc{iv}\textsuperscript{e}~siècle, tel que l’a reconstitué Ch. Pietri\footnote{Pietri, \emph{Roma christiana}, 1976 (chapitre~\ref{ch:historiographie}).} ;
\item le système de la datation consulaire et ses variations politiques : post-consulats, usurpations, consulats non reconnus\footnote{Bagnall \emph{et al.}, \emph{Consuls of the Later Roman Empire}, 1987.} ;
\item l’attraction des âges ronds et l’impossibilité d’une démographie fondée sur les âges des épitaphes\footnote{Hopkins 1966 ; Shaw 1996 ; Scheidel 2001.} ;
\item la recherche, à la fin du \textsc{iv}\textsuperscript{e}~siècle, de sépultures auprès des tombes des martyrs, et la place qu’y tient l’action de l’évêque Damase\footnote{Brown 1981 ; Diefenbach 2007 ; Trout 2015.}.
\end{itemize}

\section{Ce qui est débattu}

Quatre débats concernent directement le corpus. Le premier porte sur la valeur des épitaphes comme indicateurs des croyances. La rareté des formules doctrinales est un fait d’observation ancien, mais son interprétation reste ouverte : reflet d’une piété peu doctrinale, effet d’un genre contraint, ou décalage entre langage funéraire et langage théologique. Le deuxième porte sur le rôle de l’Église dans la gestion des cimetières et des sépultures, que les travaux de Rebillard ont remis en cause. Le troisième concerne la représentativité des inscriptions datées pour la chronologie de l’ensemble de l’épigraphie chrétienne. Le quatrième, plus récent, porte sur les conditions d’une exploitation quantitative des bases épigraphiques.

\section{Contribution du volume}

La contribution propre du volume est d’abord documentaire et méthodologique. Elle tient en six points, chacun vérifiable dans les notices, les tableaux et les données jointes.

\begin{enumerate}[label=(\arabic*),leftmargin=2em]
\item \emph{Une confrontation systématique de deux bases savantes.} \st{edbApp} notices sur \st{nA} ont été appariées à une fiche de l’Epigraphic Database Bari. \st{edbConc} textes concordent aux conventions de notation près, auxquels s’ajoutent \st{edbConcRel} textes dont la divergence apparente tient à la notation ; \st{edbExam} présentent des écarts à arbitrer ; les autres divergent fortement, le plus souvent parce que l’un des deux témoins est tronqué (chapitre~\ref{ch:transmission}).
\item \emph{Un corpus délimité par des décisions explicites.} Trois doublons sont fusionnés, deux fragments d’une même plaque réunis, une notice du \textsc{v}\textsuperscript{e}~siècle écartée. Les ensembles analytiques comptent \st{nA} notices, \st{nB} textes distincts et \st{nC} textes funéraires ; chaque mesure indique son dénominateur (chapitre~\ref{ch:corpus}).
\item \emph{Un fondement explicite pour chaque année.} Pour chaque texte, l’année est rapportée à son fondement : nom consulaire lu, lu en partie, réduit à des initiales ou restitué, formule transmise par la seule EDB, restitution conjecturale, attribution de la base. Les dix datations contestées et les cinq réserves sont publiées avec leurs motifs (chapitre~\ref{ch:fondement}).
\item \emph{Une documentation topographique et matérielle.} Les rattachements cimétériaux, les lieux de conservation, les supports, les dimensions et les signes décrits par l’EDB sont intégrés aux notices et exploités dans deux chapitres (partie~IV). Ces données étaient absentes de toutes les notices tirées de l’EDCS.
\item \emph{Une mesure du formulaire et du vocabulaire religieux sous règles publiées.} Chaque compte distingue le mot écrit, l’abréviation, la correction éditoriale, la restitution et le signe figuré (parties~V et VII).
\item \emph{Une base reproductible.} Les données consolidées, les programmes de calcul et de composition et le registre des consultations permettent de refaire chaque compte et de le contester (annexe).
\end{enumerate}

La contribution historique et théologique est plus limitée et doit être formulée avec prudence. Elle consiste à décrire le vocabulaire selon des règles contrôlables, et elle aboutit à quatre constats. Les désignations écrites de Dieu, du Christ, de l’Esprit ou de la Trinité ne concernent que \st{nProfilLarge} textes funéraires sur \st{nC}, et \st{nProfilStrict} seulement si l’on écarte les lectures qui dépendent d’une restitution ou d’un monogramme. La seule mention de la Trinité (\nnum{474}) ne comporte aucune énumération des personnes divines. Le vocabulaire de la mort est dominé par le repos et le sommeil. Enfin, la résurrection n’est nommée dans aucun texte de manière assurée : la seule attestation alléguée (\nnum{447}) est une restitution que la seconde base remplace par une autre.
